\subsection{Explicit counterexamples to the mod-$p$ good-reduction criterion}

\begin{theorem}[B-RD1: a counterexample for every odd prime]\label{thm:B-RD1}
\status{Proved.}
For every prime $p\geq3$, there is a complete discretely valued field $K$
of characteristic zero, residue field $\overline{\mathbf F}_p$, and
absolute ramification index
\[
 e(K)=\frac{p^2-1}{2},
\]
containing $\mu_p$, and an elliptic curve $E/K$ whose full $p$-torsion
is $K$-rational and whose reduction is additive, potentially good and
supersingular. Its semistability defect is exactly $2$. Its inertia
character is $-2$ on the nontrivial element of the finite inertia image.
\end{theorem}
\begin{proof}
We give the field and the curve by an explicit algebraic construction.
Put $k=\overline{\mathbf F}_p$, $R=W(k)$ and $F=\operatorname{Frac}R$.
Let $m=(p-1)/2$ and choose any root $\lambda\in k$ of
\[
 H_p(T)=\sum_{i=0}^{m}\binom mi^2T^i.
\]
Such a root exists since $m\geq1$. It is neither $0$ nor $1$:
$H_p(0)=1$ and
$H_p(1)=\binom{2m}{m}\equiv(-1)^m\pmod p$.
The coefficient calculation for Frobenius on the one-dimensional space
$H^1(\mathcal O)$ of a Legendre elliptic curve gives $(-1)^mH_p(\lambda)$
as its Hasse--Witt scalar. Thus
\[
 E_0/R:\quad y^2=x(x-1)(x-[\lambda])
\]
is smooth and has supersingular special fibre. Here $[\lambda]$ is the
Teichmuller representative. The discriminant is
$16[\lambda]^2(1-[\lambda])^2$, an $R$-unit.

Let $L=F(E_0[p])$. We now verify its ramification without invoking a
Breuil--Kisin dictionary. Every nonzero geometric $p$-torsion point
specializes to the identity, since the special fibre is supersingular.
For a local parameter $T$ of the formal group at that identity, its
height is two. Consequently $[p](T)/T$ has constant coefficient $p$,
its first coefficient which is a unit modulo $p$ is that of
$T^{p^2-1}$, and all lower nonconstant coefficients are divisible by $p$.
Weierstrass preparation gives a polynomial of degree $p^2-1$ which is
Eisenstein over $R$. Thus every nonzero $P\in E_0[p]$ generates a
completely ramified extension $F(P)/F$ of degree $p^2-1$. This degree is
prime to $p$, so the extension is tame. Since $F$ has algebraically
closed residue field, its unique tame extension of that degree is
Galois cyclic and contains every such point. Hence
\[
 \operatorname{Gal}(L/F)\simeq C_{p^2-1}.
\]
Its action on $E_0[p]\setminus\{0\}$ is free and transitive. The unique
involution in this cyclic group acts as $-1$: an involution in
$\mathrm{GL}_2(\mathbf F_p)$ with any eigenvalue $+1$ would fix a
nonzero vector.

Let $H$ be this order-two subgroup and put $K=L^H$. The Weil pairing
shows $\mu_p\subset L$, and $H$ acts on it by the determinant of $-1$,
which is $1$. Thus $\mu_p\subset K$. We have $e(K)=(p^2-1)/2$.
Choose any nonzero $P\in E_0[p]$, and put $\delta=y(P)^2\in K^\times$.
A nonzero odd-order point has $y(P)\ne0$, so $L=K(\sqrt\delta)$.
Write $E_0$ as $y^2=x^3+a_2x^2+a_4x+a_6$ and define
\[
 E:\quad y^2=x^3+\delta a_2x^2+\delta^2a_4x+\delta^3a_6.
\]
Over $L$, the map $(x,y)\mapsto(\delta x,\delta\sqrt\delta\,y)$
identifies $E_0$ with $E$. For every $Q\in E_0[p]$, both $x(Q)$ and
$\sqrt\delta\,y(Q)$ are fixed by $H$. Hence every point of $E[p]$
is $K$-rational. Equivalently, its representation is the tensor product
of the scalar quadratic representation on $E_0[p]$ and its inverse.

The curve $E_0$ has good reduction over $K$, so for every prime
$\ell\ne p$ inertia acts trivially on $T_\ell(E_0)$. On $T_\ell(E)$ it
acts by the nontrivial ramified quadratic character of $L/K$. The
Neron--Ogg--Shafarevich criterion\cite[Theorem~1]{SerreTate1968Good} therefore excludes good reduction.
It has good supersingular reduction over $L$, so its bad reduction is
additive and its minimal good-reduction inertia quotient has order $2$.
The trace on its covariant Dieudonne module is $-2$ at the involution,
since that involution acts as $[-1]$ on the special elliptic curve.
\end{proof}

\begin{theorem}[B-RD2: full $3$-torsion already over $\mathbf Q_3(\zeta_3)$]\label{thm:B-RD2}
\status{Proved, with an exact computation.}
Let $K=\mathbf Q_3(\zeta)$, where $\zeta^2+\zeta+1=0$. Then $K$ has
absolute ramification index $2$ and residue field $\mathbf F_3$. The curve
\[
 E/K:\quad y^2+y=x^3
\]
has $E[3]\subset E(K)$, additive potentially supersingular reduction,
and semistability defect exactly $6$. Its inertia character takes the
value $-2$ on the element of order $2$ of its finite inertia image.
\end{theorem}
\begin{proof}
The discriminant and $j$-invariant are $\Delta=-27$ and $j=0$.
The third division polynomial is
\[
 \psi_3(x)=3x(x^3+1).
\]
Besides the identity, the following are the eight distinct $3$-torsion
points:
\[
 (0,0),\ (0,-1),\qquad
 (x,y)\quad\text{with }x\in\{-1,-\zeta,-\zeta^2\},\quad
 y\in\{\zeta,\zeta^2\}.
\]
They exhaust $E[3]$, which has degree nine in characteristic zero.
The exact program \texttt{verify\_p3\_full\_torsion.py} checks all these
points, their orders, all $81$ sums, and that $(0,0),(-1,\zeta)$ form a
basis. Its recorded output is part of the computation certificate.

The displayed equation has $v_K(\Delta)=6$. Discriminants of two
Weierstrass equations of the same curve differ by a twelfth power.
Thus no equation over $K$ can have discriminant valuation zero, and
$E$ does not have good reduction.

We construct good reduction after ramification degree six. Make first
the unramified quadratic extension $K'=K(i)$, with $i^2=-1$, and choose
$\sqrt3\in K'$. Adjoin $r$ with $r^3=-1/4$, and then $u$ with
$u^2=r\sqrt3$. The polynomial for $s=r+1$ is
\[
 s^3-3s^2+3s-3/4.
\]
Its Newton polygon has the single slope $-2/3$ for the valuation
normalized on $K'$, so $K'(r)/K'$ is completely ramified of degree three.
The element $r\sqrt3$ has odd valuation in that field, so adjoining $u$
gives a completely ramified quadratic extension. Put $L=K'(r,u)$.
After completing the square and setting
\[
 x=r+u^2X,\qquad y+\tfrac12=u^3Y,
\]
the equation becomes
\[
 Y^2=X^3+\frac{3r}{u^2}X^2+X.
\]
The middle coefficient has positive valuation and the special fibre is
$Y^2=X^3+X$, which is smooth and supersingular in characteristic three.
Thus $E$ has good reduction over an extension of ramification degree six.

We verify minimality, including possible unramified extensions. Let
$F/K$ have good reduction. The discriminant valuation implies
$2\mid e(F/K)$. In a good integral Weierstrass equation over $F$, complete
the square; this is integral since $2$ is a unit. The resulting equation
has the form $Y^2=X^3+b_2X^2+b_4X+b_6$. Comparison with
$y^2=x^3+1/4$ uses $x=r_0+u_0^2X$, $y=u_0^3Y$, and forces
$v_3(u_0)=1/4$. Integrality of $b_6$ gives
\[
 v_3(r_0^3+1/4)\geq3/2.
\]
It follows that $r_0$ is a unit with residue $-1$. Writing $s_0=r_0+1$,
if $v_3(s_0)<1/3$ the term $s_0^3$ uniquely has lowest valuation in
$s_0^3-3s_0^2+3s_0-3/4$; if $v_3(s_0)>1/3$ the constant term uniquely
has lowest valuation. Both alternatives give valuation at most one.
Hence $v_3(s_0)=1/3$, so $3\mid e(F/\mathbf Q_3)$, and therefore
$3\mid e(F/K)$. We conclude $6\mid e(F/K)$, proving the defect is six.

Over the unramified extension $K'$, the extension just constructed is
Galois: $r\mapsto\zeta r$, $u\mapsto\zeta^2u$ has order three and
commutes with $u\mapsto-u$. Its inertia group is cyclic of order six.
The latter involution induces $Y\mapsto-Y$ on the good special fibre,
that is $[-1]$, so its trace is $-2$.
\end{proof}

\begin{proposition}[B-RD3: the ordinary case of R3(a)]\label{prop:B-RD3}
\status{Proved.}
Let $K$ be complete discretely valued of characteristic zero, with
perfect residue field of characteristic $p\geq3$, and assume
$\mu_p\subset K$. If an abelian variety $A/K$ has potentially good
ordinary reduction and $A[p]$ is unramified, then $A$ has good reduction.
There is no restriction on $\dim A$.
\end{proposition}
\begin{proof}
We may pass to the completed maximal unramified extension: good
reduction is equivalent before and after this passage by the
Neron--Ogg--Shafarevich criterion. The residue field is now algebraically
closed. Choose a finite Galois extension $L/K$ giving good reduction,
and let $\mathcal A/\mathcal O_L$ be the resulting abelian scheme.
For its ordinary special fibre, the connected subgroup
$\mathcal A[p]^0$ is of multiplicative type and its etale quotient has
rank $p^g$. Indeed the Cartier dual of the connected subgroup has
etale special fibre, hence is etale by Nakayama's lemma applied to its
relative differentials. Over the strictly henselian ring
$\mathcal O_L$, it is constant; thus
$\mathcal A[p]^0\simeq\mu_p^g$.

Every inertia element acts semilinearly on this subgroup by its action
on the base, composed with a matrix in $\mathrm{GL}_g(\mathbf F_p)$.
The base action fixes $\mu_p$, and its generic action on $A[p]$ is
trivial. The matrix is therefore the identity. The action on the etale
quotient is trivial for the same reason. Over the algebraically closed
residue field the connected--etale sequence splits, and there are no
nonzero homomorphisms between the connected and etale factors in either
direction. Thus inertia acts identically on $\mathcal A_k[p]$.

If $\alpha$ is one of the resulting finite-order automorphisms of
$\mathcal A_k$, then $\alpha-1$ factors through $[p]$, so
$\alpha=1+p\beta$ for some $\beta\in\operatorname{End}(\mathcal A_k)$.
The kernel of reduction modulo $p$ in the units of this endomorphism
ring has no nontrivial finite-order elements for $p\geq3$. To see this,
pass to its finite free $\mathbf Z_p$-completion. If
$\alpha=1+p^a\gamma$ with $a\geq1$ and $\gamma\notin p\operatorname{End}$,
the binomial expansion shows that a prime-to-$p$ power has the same
leading valuation, and the $p$-th power has leading valuation $a+1$;
neither can be one. This proves the assertion for any finite order.
Hence inertia is trivial on the entire special abelian variety.
Proper smooth base change now makes $T_\ell(A)$ unramified for
$\ell\ne p$, and the same good-reduction criterion concludes.
\end{proof}

\begin{corollary}[B-RD6: the manuscript's cas test is negative]\label{cor:B-RD6}
\status{Refuted, for every odd prime.}
For every $p\geq3$ there is a field $K$ as in Theorem~\ref{thm:B-RD1}
and two elliptic curves $E,E'/K$ with potentially good reduction and
isomorphic $p$-torsion, for which the inertia traces differ on an element
of order prime to $p$.
\end{corollary}
\begin{proof}
Use $E$ and $K$ from Theorem~\ref{thm:B-RD1}. Choose an ordinary
elliptic curve $C/\overline{\mathbf F}_p$; for example choose a Legendre
parameter outside the finite root set of $H_p$ and outside $\{0,1\}$.
The ordinary Serre--Tate coordinates of
\cite[Theorem~2.1]{Katz1981SerreTate} give its canonical deformation
over $W(\overline{\mathbf F}_p)$, whose $p$-divisible group is
\[
 \mu_{p^\infty}\oplus\mathbf Q_p/\mathbf Z_p.
\]
This formal elliptic scheme algebraizes: its identity section provides
the compatible relatively ample bundles $\mathcal O(3[0])$, so the
formal existence theorem applies. Let $E'$ be its generic fibre after
extension to $K$. Then $E'$ has good reduction, and $\mu_p\subset K$
gives $E'[p]\simeq(\mathbf Z/p)^2\simeq E[p]$ over $K$.
The good curve has inertia trace $2$; the bad quadratic twist has trace
$-2$ at the inertia involution. That involution has order prime to $p$.
\end{proof}
